(1) 从平面镜出发往返一次，等效矩阵可取 \(M=PSP\)，平面镜反射不引入曲率矩阵。连乘得 \(A=D=1-2L/R,\ B=2L(1-L/R),\ C=-2/R\)，行列式为1。周期线性光学映射稳定要求 \(|(A+D)/2|<1\)，即 \(|1-2L/R|<1\)，因此稳定区是 \(0<L<R\)。边界 \(L=0\) 或 \(L=R\) 时迹等于 \(\pm2\)，不属于严格束缚区。

(2) 腔的唯一轴对称基模必须在平面镜处具有平面相位前沿，即束腰。沿自由传播距离 \(z\)，高斯束曲率满足 \(R_b(z)=z[1+(z_R/z)^2]\)。在曲面镜 \(z=L\) 处应有 \(R_b(L)=R\)，于是 \(R=L+z_R^2/L\)，即 \(z_R^2=L(R-L)\)，从而 \(z_R=\sqrt{L(R-L)}\)。只有 \(0<L<R\) 才存在实且正的束腰瑞利长度，这独立验证矩阵稳定条件。

(3) 按 \(z_R=\pi w_0^2/\lambda\)，束腰半径平方 \(w_0^2=(\lambda/\pi)\sqrt{L(R-L)}\)。凹镜处束斑 \(w_L^2=w_0^2[1+(L/z_R)^2]=(\lambda/\pi)R\sqrt{L/(R-L)}\)。两式量纲均为长度平方；当 \(L=R/2\) 时 \(z_R=R/2\)，可直接代回波前曲率方程核验。

(4) 若 \(L\to0^+\)，\(w_0\to0\)，近轴高斯模式可能受到衍射、有限光阑和镜面缺陷限制，理想腔模型不保证无限小束腰可实现。若 \(L\to R^-\)，\(w_L\to\infty\)，有限镜口会产生严重裁剪损耗，尽管数学上仍处于严格稳定区内。实际激光腔设计应同时检查镜口、衍射损耗和热透镜等因素，不能仅凭 \(|\operatorname{tr}M|<2\) 保证低损耗。
